\documentclass[11pt,a4paper]{article}
\usepackage{amsmath,amssymb,amsthm}
\usepackage[margin=2.5cm]{geometry}
\usepackage{hyperref}

\newtheorem{definition}{Definition}[section]
\newtheorem{theorem}[definition]{Theorem}
\newtheorem{proposition}[definition]{Proposition}
\newtheorem{remark}[definition]{Remark}
\newtheorem{conjecture}[definition]{Conjecture}

\title{Hyperseed Formalization:\\
  Robust Cognitive Strategies for Resource Abundance}
\author{ProtomegaTron (automated formalization)}
\date{2026-09-18}

\begin{document}
\maketitle

\begin{abstract}
We formalize Ben Goertzel's ``Robust Cognitive Strategies for Resource
Abundance'' (2023-04-26) within the Hyperseed ontology. Key mappings:
resource scarcity as thin fiber constraint, abundance as thick fiber
enablement, heuristics as fiber shortcuts, the inability to foresee
abundance strategies as a fiber horizon, and attention as fiber allocation.
\end{abstract}

\tableofcontents

\section{Fiber thickness: scarcity vs abundance}

\begin{theorem}[Resource regime determines fiber thickness --- claims 1--10, 19--20]
Resource scarcity forces thin fiber --- limited perspectives, limited
inference paths, because each path costs compute:
\[
  \text{Resources} \downarrow \implies \text{thickness}(F) \downarrow
  \implies |\text{perspectives}| \downarrow
\]
Resource abundance enables thick fiber --- many perspectives, exhaustive
inference, ensemble reasoning:
\[
  \text{Resources} \uparrow \implies \text{thickness}(F) \uparrow
  \implies \text{qualitatively different strategies}
\]
The transition from thin to thick fiber is qualitative, not quantitative.
Thick fiber enables cognitive strategies (brute-force search, massive
parallelism, exhaustive verification, redundancy) that are structurally
impossible under thin fiber, not just expensive.
\end{theorem}

\section{Fiber shortcuts: heuristics under scarcity}

\begin{proposition}[Heuristics as compressed fiber --- claims 2, 11, 21]
Heuristics are fiber shortcuts --- compressed inference paths that trade
accuracy for speed:
\[
  h: F_{\text{full}} \to F_{\text{compressed}} \quad \text{with} \quad
  \text{cost}(F_{\text{compressed}}) \ll \text{cost}(F_{\text{full}})
\]
These shortcuts become ingrained under scarcity and feel like ``the right
way to think.'' Under abundance, full inference becomes feasible and
shortcuts become suboptimal --- but systems shaped by scarcity continue
to use them from habit.
\end{proposition}

\section{Fiber horizon: can't foresee thick from thin}

\begin{theorem}[Fiber horizon --- claims 3, 11--14, 22]
The inability to foresee abundance strategies from within scarcity is a
\emph{fiber horizon}: thin fiber literally can't represent what thick
fiber would do:
\[
  F_{\text{thin}} \not\models \text{description}(F_{\text{thick}})
\]
Thick-fiber solutions are invisible from within thin fiber --- they
``feel wasteful'' because thin fiber has no way to evaluate their
benefits. This is not ignorance but a structural limitation: the
representational capacity of thin fiber is insufficient to model
thick-fiber reasoning.
\end{theorem}

\section{Fiber allocation: attention under scarcity}

\begin{proposition}[Attention as fiber allocation --- claims 18, 23--24]
Attention (selecting what to think about) is fiber allocation under
scarcity --- choosing which fibers to maintain given limited resources:
\[
  \text{Attention}(b) = \arg\max_S \text{value}(S) \quad
  \text{subject to} \quad |S| \leq \text{budget}
\]
Under abundance, attention becomes less constrained. Ensemble methods
(maintaining multiple fiber sections simultaneously) become feasible ---
this is multi-section fiber, comparing outputs from multiple perspectives
rather than committing to one.
\end{proposition}

\section{Cross-references}

\begin{remark}[Connections]
\begin{itemize}
\item \textbf{A BGI Manifesto} (2023-11-25): intelligence explosion
  as transition to compute abundance, enabling qualitatively different
  cognitive strategies.
\item \textbf{Seeding RSI} (2026-07-28): self-improvement under abundance
  --- exploring many improvement paths simultaneously.
\item \textbf{GPT o1 Does Not Know What It's Doing} (2024-10-03):
  current LLMs as scarcity-shaped cognition --- pattern matching as
  a fiber shortcut.
\item \textbf{Introducing Hyperseed-1} (2024-11-27): ontology as
  inference accelerator --- a scarcity-regime strategy for speeding
  up reasoning with limited resources.
\end{itemize}
\end{remark}

\end{document}
